\chapter{The model}\label{ch:model}

\begin{chapterplan}
The question of this chapter: what are the rules of the game, who knows what at each date, and what must an equilibrium of the game satisfy?
After this chapter you can:
\begin{itemize}
\item Name every player and say what each one observes.
\item Derive the order probabilities, the no-entry proceeds \(t_0\), the entry probability, and high-class ownership from the rules.
\item State the six equilibrium conditions in words and derive the investor's objective.
\item Say why a unilateral deviation differs from a comparison across economies.
\item Explain why a common rescale of money changes no order, belief, entry, or ownership.
\end{itemize}
\end{chapterplan}

This chapter sets out the primitives, the rules, and the definition of equilibrium. It contains no result of the paper.
Read it slowly, because a misread assumption changes what a later result says.


%==================================================================
\section{The decision interval}\label{sec:model-interval}

A listed company is for sale, and the sale is public.
One buyer already did the preparation work. A second buyer must still decide whether to do it, and the work costs money.
While the second buyer decides, the stock of the target still trades.
An outside investor knows how well the second buyer fits the target, and this investor trades the stock.
The second buyer reads the price.

\subsection*{The institutional set-up}

The model applies to a listed target with a public sale opportunity.
The stock must still trade, and a further bidder must still be able to join.
For example, a disclosed approach, an announced strategic review, or an open contest can create such an interval.
In each case the opportunity must become public before the set of buyers is fixed.
Only then can a prospective buyer see the price during the period that matters.
Chapter~\ref{ch:empirical} explains how to date this interval in a given transaction.

Two buyers appear in the interval.
The first is the incumbent. It is the party that approached the target, or a bidder that completed diligence.
The public can infer its strength from the disclosed record.
The second is the challenger. It must still decide whether to prepare.

\subsection*{Why preparation is costly and sunk}

A buyer cannot submit an executable proposal before it prepares.
Preparation has three parts:
\begin{itemize}
\item verify the target's contracts, liabilities, and technology;
\item arrange funds and approvals;
\item commit evaluation capacity that has other uses.
\end{itemize}
The buyer cannot recover these costs after it pays them. They are sunk.
They differ across buyers.
Prior familiarity, spare capacity, or a small deal relative to the buyer's size make preparation cheap for some buyers.
For other buyers the same work is expensive.

A challenger that does not prepare stays out of the sale.
Its cost is therefore an entry cost. Payment of the cost reveals the challenger's value and permits a bid.

\subsection*{Why the investor and the challenger know different things}

The challenger knows its own integration plans.
A specialist investor can know what the target's customers do, how its technology compares with alternatives, or how demand for its product will change.
These are different facts about the same acquisition match.
A more accurate signal about one fact can leave the other fact unknown.

The model takes a simple form of this idea.
The investor knows the challenger's quality class exactly. It does not know the challenger's value.
The challenger starts with no information of its own.
Chapter~\ref{ch:robustness} removes both extremes, and Chapter~\ref{ch:research} makes the investor's information the result of costly research.
The identity of the challenger is common knowledge. The market is uncertain only about how well the challenger fits the target.

\subsection*{Incumbent strength is a distribution}

The public sees a record, and the record does not give a bid.
It can show that a strategic acquirer with clear synergies is at the table.
It can show that the prepared party is a financial buyer with a leveraged capital structure.
In both cases the record can stay silent on what the buyer would pay.
Strategic and financial bidders differ systematically in their valuations (Gorbenko and Malenko, 2014, cited in the paper).

The model therefore treats incumbent strength as a distribution.
It varies that distribution, and it keeps the incumbent's realized value private until the incumbent bids.
For the preparation stage, the paper cites the definitive proxy statement of Imprivata (2016).
That document separates an unsolicited approach, outreach to potential buyers, and indications of interest that depend on further diligence.
The paper also cites studies in which takeover competition develops before public bids (Boone and Mulherin, 2007; Gentry and Stroup, 2019).

\begin{analogy}[an art sale]
An art dealer must pay an expert to examine a painting before the dealer may bid.
Before the dealer pays, it can look at the market price of a claim on the sale.
That price moves with the trades of an outside collector who knows how well the painting fits the dealer's collection.
Other people trade the claim for reasons that have no link to the painting.
So the dealer reads the price as a noisy hint about its own fit.

Where the analogy breaks: an art sale has no traded claim on its proceeds.
In the model the claim is a target share, and its payoff depends on whether the challenger prepares.
So the price also reflects the challenger's own response to the price.
\end{analogy}

{\small\color{notegrey}Source: Section~2.1 of the paper.}

%==================================================================
\section{The players}\label{sec:model-players}

The model has six kinds of agent. All strategic agents are risk neutral.
The target share is the only traded claim with information about the acquisition match, and neither bidder trades it.

The target is the listed company, and its board is the seller.
Before any trade, the seller commits in public to a cash second-price auction with reserve \(p>0\).
After this announcement the seller takes no further action in the model.
The sale binds all target shares, so the tender decisions and holdout of shareholders are outside the model.

The incumbent is the prepared buyer. It pays no further entry cost.
It knows its acquisition value \(R\) before it bids.
The market knows only the distribution \(F\) of \(R\), and \(F\) is public when the sale opportunity becomes visible.

The challenger is not yet prepared. At the start it knows neither its class nor its value.
It sees the stock price. Then it learns its preparation cost \(C\) in private.
Then it decides whether to prepare. If it pays \(C\), it learns its value and may bid. If it does not pay, it stays out of the sale.

The informed investor is an outside sector specialist, such as a sector-focused long-short fund.
It holds no net position in the target before trade, has no control rights, and cannot acquire the target.
It forms its view from its own research on the target's customers, technology, and demand.
It sees the challenger's class \(j\). A nonzero order costs it \(k>0\).

The liquidity trader trades for reasons that have no link to the sale.
A liquidity sale comes from a current shareholder who sells for cash. A liquidity purchase comes from an outside investor.
An abstention means that the order opportunity passes with no order.

The market makers are competitive and risk neutral, and they earn zero expected profit.
They see the order and nothing else.
They set the price from the order, and they absorb the order at that price.
Each share held at the sale, market-maker inventory included, receives the terminal payoff \(V_T\).

{\small\color{notegrey}Source: Sections~2.2 and 2.3 of the paper.}

%==================================================================
\section{Values and classes}\label{sec:model-values}

Two values decide how the auction ends: what the incumbent would pay, and what the challenger would pay.
The incumbent's value is a draw from a known distribution.
The challenger's value is a draw from a distribution that depends on its class.
Class \(H\) is the class with the better fit, and class \(L\) is the class with the worse fit.
The target's known standalone value is normalized to zero, and all values and costs are measured per target share.

\subsection*{The incumbent's value}

The incumbent's value \(R\) has a continuous distribution \(F\) on \([0,\bar r]\).
Here \(F(y)=\Pr(R\le y)\), and \(F\) is extended by unity above \(\bar r\): \(F(y)=1\) for \(y\ge\bar r\).
The extension lets you write \(F(\theta)\) for every \(\theta\ge0\), also for a \(\theta\) above \(\bar r\). The incumbent's value is below such a \(\theta\) for sure.

Continuity has a clear use.
A continuous distribution function has no jump, so \(\Pr(R=y)=F(y)-F(y^-)=0\) at every \(y\).
The incumbent's value equals the reserve with probability zero.
For a challenger value \(\theta\), a tie \(R=\theta\) also has probability zero. For each fixed \(\theta\), the independent value \(R\) hits that one point with probability zero.
The auction therefore needs no rule for equal bids from the two buyers.

\subsection*{Strength, strengthening, and a path}

A stronger incumbent has a value that is likely to be higher.
The model makes this exact with first-order stochastic dominance (Section~\ref{sec:toolkit-fosd}).
A distribution \(\tilde F\) is a strengthening of \(F\) when \(\tilde F(y)\le F(y)\) for every \(y\).
In words: at every level \(y\), the stronger incumbent is less likely to have a value at or below \(y\).

A path is a family \(\{F_r\}\) of continuous incumbent distributions indexed by a strength \(r\) in an interval.
Along the path, \(F_{r'}\) is a strengthening of \(F_r\) whenever \(r'>r\).
Everything else stays fixed, for example the classes, the prior, the reserve, the costs, and the slot.
The example of Chapter~\ref{ch:example} uses such a path: \(F_r\) is uniform on \([0,r]\), so \(F_r(y)=y/r\) for \(0\le y\le r\).
A larger \(r\) gives a smaller \(F_r(y)=y/r\) at each level \(y\) with \(0<y<r\).

\subsection*{The challenger's class and value}

The challenger belongs to one of two classes, \(j\in\{H,L\}\).
The prior is \(\pi=\Pr(j=H)\), and \(0<\pi<1\).
Given its class \(j\), the challenger's value \(\theta\) has a distribution \(G_j\) on \([0,\infty)\) with a finite mean.
Atoms are allowed. The example in Chapter~\ref{ch:example} uses a single point mass in each class.

Two pieces of notation come with \(G_j\).
The number \(G_j(p^-)\) is the probability that a class-\(j\) value falls below the reserve \(p\).
The function \(D(y)=G_L(y)-G_H(y)\) is the difference between the two class distributions at \(y\).

\subsection*{The class ranking}

The paper assumes two restrictions on the classes:
\begin{equation}\label{eq:model-classes}
G_H(y)\le G_L(y)\quad\text{for every }y,\qquad G_H(p^-)=G_L(p^-)=0.
\end{equation}
The paper calls the two restrictions in \eqref{eq:model-classes} the class ranking.

Read the two halves one at a time.
\begin{itemize}
\item The first half says that class \(H\) dominates class \(L\). At every level \(y\), a class-\(H\) value is no more likely than a class-\(L\) value to lie at or below \(y\). In symbols, \(D(y)\ge0\) for every \(y\).
\item The second half says that no value of either class lies below the reserve. Every value of either class meets the reserve, so the reserve excludes no class. In symbols, \(D(p^-)=G_L(p^-)-G_H(p^-)=0\).
\end{itemize}

\begin{watchout}
The two halves of \eqref{eq:model-classes} do different jobs.
The first half says which class is better. The second half says that the reserve treats the classes alike.
In Chapter~\ref{ch:auction} the reserve adds a term \(pF(p)D(p^-)\) to the spread, and the second half sets that term to zero.

The paper finds the weakest restriction of this kind (Proposition~1(c) in Appendix A.1).
Suppose that neither class has a value atom at the reserve.
Then the spread weakly rises under every strengthening of every continuous \(F\) exactly when two conditions hold:
\begin{itemize}
\item the high class dominates above the reserve, \(D(y)\ge0\) for every \(y>p\);
\item the two classes fall below the reserve with equal probability, \(G_H(p^-)=G_L(p^-)\).
\end{itemize}
The paper calls the second condition equal admissibility.
Chapter~\ref{ch:welfare} studies a reserve that excludes the low class, where equal admissibility fails.
\end{watchout}

The paper places this structure in the selective-entry literature on auctions.
The class indexes the distribution of the value, preparation reveals the value, and the stock price is the signal that comes before entry.

{\small\color{notegrey}Source: Section~2.2 of the paper.}

%==================================================================
\section{Preparation costs}\label{sec:model-costs}

After the challenger sees the stock price, it privately learns its preparation cost \(C\in\{0,c_H\}\).
The cost is zero with probability \(\rho\), so \(\Pr(C=0)=\rho\) and \(0<\rho<1\). Otherwise the cost is \(c_H>0\).
A zero-cost challenger is a ready buyer: its diligence costs nothing at the margin.
A high-cost challenger pays \(c_H\).
In both cases a payment of \(C\) reveals \(\theta\) and permits a bid, and a challenger that does not pay stays outside the sale.

\subsection*{A zero-cost challenger always prepares}

If the challenger prepares and bids its value, its gross profit is an expected auction rent.
Let \(Y\) be the highest competing admissible bid, with \(Y=0\) if there is none.
A bidder with value \(\theta\ge p\) that bids its value earns \((\theta-\max\{p,Y\})_+\), where \((x)_+=\max\{x,0\}\) is the positive part (Section~\ref{sec:toolkit-tail}).
A bidder with \(\theta<p\) does not bid and earns zero. In both cases the amount is never negative.
Its expectation is therefore never negative either.

A zero-cost challenger pays nothing for this rent, so preparation cannot hurt it.
The definition of equilibrium fixes the choice: a zero-cost challenger prepares after every price.
Two consequences follow later in this chapter: entry never falls below \(\rho\), and the choice of a zero-cost challenger does not depend on its belief.

\subsection*{The independence assumptions}

The paper states that four objects are mutually independent:
\begin{itemize}
\item the pair \((j,\theta)\), that is, the class and the value;
\item the incumbent's value \(R\);
\item the preparation cost \(C\);
\item the draws of the order slot, that is, the draw that selects who fills the slot and the draw of the liquidity order.
\end{itemize}
The paper draws two consequences.
First, the cost carries no information about the class or the value.
Second, given the class, the challenger's value is independent of the order, because the order depends on the class and the slot draws only.

Each assumption has a job. The list below is the guide's reading, and the paper names the uses in the sections cited.
\begin{itemize}
\item \emph{The cost is independent of the class.} The challenger learns \(C\) after it sees the price. If \(C\) were informative about the class, the challenger would revise its belief when it learns \(C\). Then the belief at the price would be the wrong input. Entry would also differ by class, and the expected payoff of Section~\ref{sec:model-equilibrium} would need an entry probability for each class.
\item \emph{The incumbent's value is independent of everything else.} The price carries no news about \(R\). The numbers \(t_0\), \(t_j\), and \(g_j\) then depend on \(F\) and the reserve, and on no belief.
\item \emph{The slot draws are independent of the class.} Who fills the slot is not a signal of the class. Only the investor's strategy connects the order to the class.
\item \emph{The value is independent of the order, given the class.} The investor knows the class but not the value. Its order therefore says no more about the value than the class says. This is the reason why the challenger's expected profit is linear in its belief (Section~\ref{sec:model-payoffs}).
\end{itemize}

{\small\color{notegrey}Source: Sections~2.2 and 2.3 of the paper, and Definition A.1.}

%==================================================================
\section{The sale}\label{sec:model-sale}

Before trade, the seller commits to a cash second-price auction with reserve \(p>0\).
The rules are as follows.
\begin{itemize}
\item Only a prepared bidder can bid.
\item A bid is admissible when it is at least \(p\). A bid equal to the reserve is admissible.
\item The highest admissible bidder acquires the target.
\item The winner pays the larger of the reserve and the highest competing admissible bid.
\item If there is no admissible bid, there is no sale.
\end{itemize}

\subsection*{Bids equal values}

Each prepared bidder whose value meets the reserve bids its value.
The bid is weakly dominant. Chapter~\ref{ch:auction} proves this, and the idea is short.

In a second-price auction, your bid decides whether you win. It does not decide what you pay.
A bid above your value can win at a price above your value, when a rival bids above your value.
A bid below your value can lose an auction that you want to win.
A prepared challenger whose value is below the reserve does not bid.
Under the class ranking this case does not occur.

\subsection*{The proceeds without entry}

Suppose the challenger does not prepare. Then only the incumbent can bid.
The next derivation finds the expected amount that the target receives in this case.
The target's standalone value is zero, so the proceeds are zero when there is no sale.

\begin{derivation}[the no-entry proceeds \(t_0\)]
\begin{enumerate}[leftmargin=2em]
\item The incumbent bids its value \(R\), because bids equal values.
\item The bid is admissible exactly when \(R\ge p\).
\item If the bid is admissible, no rival bid exists. The highest rival bid is zero, and the incumbent pays \(\max\{p,0\}=p\).
\item If the bid is not admissible, there is no sale, and the target receives zero.
\end{enumerate}
Steps 1 to 4 give the target's proceeds, \(p\,\1\{R\ge p\}\), where \(\1\{A\}\) is one when \(A\) holds and zero otherwise.
Take the expectation:
\begin{align*}
t_0&=\E\bigl[p\,\1\{R\ge p\}\bigr] \why{proceeds without entry}\\
&=p\,\E\bigl[\1\{R\ge p\}\bigr] \why{$p$ is a constant}\\
&=p\,\Pr(R\ge p) \why{the mean of an indicator is a probability}\\
&=p\,\bigl[1-\Pr(R<p)\bigr] \why{complement rule}\\
&=p\,\bigl[1-F(p)\bigr] \why{$F$ is continuous, so $\Pr(R<p)=F(p)$}
\end{align*}
The last step uses the left limit of the distribution function (Section~\ref{sec:toolkit-cdf}): in general \(\Pr(R<p)=F(p^-)\), and continuity of \(F\) gives \(F(p^-)=F(p)\).
\end{derivation}

The result is
\begin{equation}\label{eq:model-t0}
t_0=p\,[1-F(p)].
\end{equation}

In words: the board collects its announced walk-away price \(p\) from a lone prepared bidder, and the bidder pays it only when its value meets \(p\).
The factor \(1-F(p)\) is the chance of that event.

Test the formula at the extremes.
If the reserve is above \(\bar r\), then \(F(p)=1\), the incumbent never meets the reserve, and \(t_0=0\).
If the reserve is near zero, then \(t_0\le p\) is near zero, because the incumbent pays \(p\) when it buys.
In between, \(t_0\) is positive whenever \(F(p)<1\), that is, whenever the incumbent meets the reserve with positive probability.
In the example of Chapter~\ref{ch:example}, \(F(p)=p/r\), so \(t_0=p(1-p/r)\) (Appendix A.1 of the paper).
That expression is zero at \(p=0\) and at \(p=r\), and it is positive between.

\begin{checkyourself}
Find \(t_0\) when \(p\ge\bar r\).
Show that \(0\le t_0\le p\) for every reserve \(p>0\).
In the example, with \(F(p)=p/r\) for \(0<p<r\), write \(t_0\) as a product and say where it is zero.
\label{q:model-t0}
\end{checkyourself}

{\small\color{notegrey}Source: Section~2.2 of the paper and Appendix A.1.}

%==================================================================
\section{The order slot}\label{sec:model-slot}

Trade takes place in one order slot.
The slot is a single public order opportunity, and it comes before the challenger decides.
Exactly one trader fills it. The market makers see the order and not the trader.

A purchase can therefore come from the investor, who has good news, or from a liquidity buyer, who has none.
A purchase is a hint about the class, and it is never proof.
A picture to keep in mind: the clerk of a letter box finds one slip and cannot tell which of two writers wrote it.
The picture breaks at an empty box. In the slot, no order is also an observation, the order \(0\), and it can carry news about the class too.

\subsection*{The rules of the slot}

With probability \(\lambda\in(0,1)\), the informed investor fills the slot.
It sees the class \(j\) and chooses an order \(q\in\{-1,0,1\}\).
The order \(-1\) is a sale of one unit, \(0\) is an abstention, and \(1\) is a purchase of one unit.
A nonzero order costs \(k>0\).

With probability \(1-\lambda\), a liquidity trader fills the slot instead.
It sells, abstains, or buys one unit with probabilities \(\beta\), \(1-2\beta\), and \(\beta\).
Write \(\eta\) for this liquidity order distribution:
\[
\eta(-1)=\eta(1)=\beta,\qquad \eta(0)=1-2\beta.
\]
The paper requires \(0<\beta\le1/3\) in the definition of the slot. The end of this section explains this domain.
No one observes who filled the slot, and the identity draw is independent of the class, the costs, and the incumbent's value.

\subsection*{The probability of each order}

Let \(X\in\{-1,0,1\}\) be the order that the market makers see.
Fix a class \(j\).
The investor's strategy in class \(j\) is a probability distribution \(\sigma_j=\bigl(\sigma_j(-1),\sigma_j(0),\sigma_j(1)\bigr)\) on the three orders.
It can be pure or mixed, and a wrong-signed order is allowed.
Write \(f_j(x)\) for the probability of the order \(x\) in class \(j\).
Figure~\ref{fig:model-tree} shows the two ways to reach each order.

\begin{figure}[tbp]
\centering
\begin{tikzpicture}[font=\small,>=Stealth,every node/.style={inner sep=2pt}]
\node[anchor=west,text=notegrey] at (-0.3,3.4) {who fills the slot};
\node[anchor=west,text=notegrey] at (6.3,3.4) {order};
\node[anchor=west,text=notegrey] at (9.0,3.4) {probability of the path};
\node (root) at (0.2,0) {slot};
\node (inv) at (2.7,1.7) {investor};
\node (liq) at (2.9,-1.7) {liquidity trader};
\draw (root.east) -- node[above,sloped,pos=0.5] {\(\lambda\)} (inv.west);
\draw (root.east) -- node[below,sloped,pos=0.5] {\(1-\lambda\)} (liq.west);
\foreach \y/\q/\s/\pp in {2.6/-1/{\sigma_j(-1)}/{\lambda\,\sigma_j(-1)},1.7/0/{\sigma_j(0)}/{\lambda\,\sigma_j(0)},0.8/1/{\sigma_j(1)}/{\lambda\,\sigma_j(1)}}{
  \node[anchor=west] (i\q) at (6.3,\y) {\(X=\q\)};
  \draw (inv.east) -- node[above,pos=0.78,fill=white,inner sep=1pt] {\(\s\)} (i\q.west);
  \node[anchor=west] at (9.0,\y) {\(\pp\)};
}
\foreach \y/\q/\s/\pp in {-0.8/-1/{\beta}/{(1-\lambda)\beta=b},-1.7/0/{1-2\beta}/{(1-\lambda)(1-2\beta)=b_0},-2.6/1/{\beta}/{(1-\lambda)\beta=b}}{
  \node[anchor=west] (l\q) at (6.3,\y) {\(X=\q\)};
  \draw (liq.east) -- node[above,pos=0.78,fill=white,inner sep=1pt] {\(\s\)} (l\q.west);
  \node[anchor=west] at (9.0,\y) {\(\pp\)};
}
\end{tikzpicture}
\caption{The probability tree of the order slot in class \(j\). The probability of an order is the sum of the two path probabilities that end in it.}\label{fig:model-tree}
\end{figure}

\begin{derivation}[the probability \(f_j(x)\) of an order]
Let \(I\) be the event that the investor fills the slot.
Then \(\Pr(I\mid j)=\lambda\), because the identity draw is independent of the class.
The events \(I\) and its complement partition the possible outcomes, so total probability applies (Section~\ref{sec:toolkit-bayes}):
\begin{align*}
f_j(x)&=\Pr(X=x\mid j)\why{definition of $f_j$}\\
&=\Pr(I\mid j)\Pr(X=x\mid I,j)+\Pr(I^c\mid j)\Pr(X=x\mid I^c,j)\why{total probability}\\
&=\lambda\,\sigma_j(x)+(1-\lambda)\,\eta(x)\why{investor: $\sigma_j$; liquidity trader: $\eta$}
\end{align*}
Check that the three probabilities add to one:
\begin{align*}
\sum_x f_j(x)&=\lambda\sum_x\sigma_j(x)+(1-\lambda)\sum_x\eta(x)\why{sum each term over $x$}\\
&=\lambda\cdot1+(1-\lambda)\bigl[\beta+(1-2\beta)+\beta\bigr]\why{both are probability distributions}\\
&=\lambda+(1-\lambda)=1\why{$\beta+1-2\beta+\beta=1$}
\end{align*}
\end{derivation}

Next, bound the probability of a nonzero order.
Define
\[
b=(1-\lambda)\beta.
\]
A liquidity order in a given direction arrives with probability \(b\).
For \(x\in\{-1,1\}\), the liquidity term is \((1-\lambda)\eta(x)=(1-\lambda)\beta=b\).
Since \(0\le\sigma_j(x)\le1\) and \(\lambda>0\), the informed term obeys \(0\le\lambda\sigma_j(x)\le\lambda\).
Add \(b\) to each side of this double inequality:
\begin{equation}\label{eq:model-orderprob}
f_j(x)=(1-\lambda)\,\eta(x)+\lambda\,\sigma_j(x),\qquad b\le f_j(x)\le b+\lambda\quad\text{for }x\in\{-1,1\}.
\end{equation}

In words: the lower end occurs when the investor never places the order \(x\). Then only a liquidity trader places it, with probability \(b\).
The upper end occurs when the investor always places \(x\). Then the order arrives with probability \(b+\lambda\).
The zero order works in the same way, with \(b_0=(1-\lambda)(1-2\beta)\) in place of \(b\): \(b_0\le f_j(0)\le b_0+\lambda\).

Test the extremes.
If \(\lambda\) is near zero, the investor is almost never present.
Then \(f_j(x)\) is near \((1-\lambda)\eta(x)\) in both classes, and the order carries almost no news.
If \(\beta\) is small, liquidity trade in each direction is rare, so \(b\) is small and one order from the investor is a strong signal.

\subsection*{The largest likelihood ratio of one order}

Define
\[
u=\frac{\lambda}{b},
\]
the informed presence mass relative to the liquidity mass in one direction.

The likelihood ratio of an order \(x\) is \(f_H(x)/f_L(x)\).
By \eqref{eq:model-orderprob}, for a nonzero order, the numerator is at most \(b+\lambda\) and the denominator is at least \(b\).
Both are positive. The ratio is therefore at most
\[
\frac{b+\lambda}{b}=1+\frac{\lambda}{b}=1+u.
\]
The same argument with the roles of the classes swapped gives a lower bound \(1/(1+u)\).
The upper bound is reached when class \(H\) always places \(x\) and class \(L\) never does.
So \(1+u\) is the largest likelihood ratio that one order can carry.

This bound limits what the market makers and the challenger can learn.
By Bayes' rule in odds form (Section~\ref{sec:toolkit-odds}), the posterior odds equal the prior odds times the likelihood ratio.
Chapter~\ref{ch:beliefs} proves that, under every investor strategy, the posterior after every order lies between a bottom posterior \(m(\pi)\) and a top posterior \(M(\pi)\).
For the zero order this uses \(\beta\le1/3\), as the end of this section explains.
The two bounds are
\begin{equation}\label{eq:model-bounds}
M(\pi)=\frac{\pi(1+u)}{\pi(1+u)+1-\pi},\qquad
m(\pi)=\frac{\pi}{\pi+(1-\pi)(1+u)}.
\end{equation}
This chapter only states the bounds. They rest on one property of the slot: each class reaches each order through a liquidity trader, so the investor keeps a residual advantage after every order.

\subsection*{The domain of \texorpdfstring{\(\beta\)}{beta}}

The paper defines the slot with \(0<\beta\le1/3\).
A footnote in Section~2.3 of the paper, and a paragraph on the liquidity share in Appendix A.2, say which results use the bound \(\beta\le1/3\).
In plain words:
\begin{itemize}
\item The posterior bound common to all orders uses \(\beta\le1/3\). So does everything that rests on it: the exclusion bound (Corollary~A.1, Chapter~\ref{ch:beliefs}) and the part of the research bounds that concerns the zero order.
\item The parameter region of the research chapter, and the identification statements on which slots are feasible (Appendix A.8), consider only the slots of the model. These are the slots with \(\beta\le1/3\).
\item Every other result holds for every \(0<\beta<1/2\). A corollary that fixes a particular slot holds at that slot.
\end{itemize}

Two facts explain the two numbers.
The first is \(\beta<1/2\). The probability \(\eta(0)=1-2\beta\) must be positive.
Then \(f_j(0)\ge(1-\lambda)\eta(0)>0\), and every order has positive probability in both classes.
The second is \(\beta\le1/3\). The zero order has a liquidity term \(b_0=(1-\lambda)(1-2\beta)\). Compare it with \(b\):
\begin{align*}
b_0\ge b&\iff(1-\lambda)(1-2\beta)\ge(1-\lambda)\beta\why{definitions of $b_0$ and $b$}\\
&\iff1-2\beta\ge\beta\why{divide by $1-\lambda>0$}\\
&\iff1\ge3\beta\why{add $2\beta$ to both sides}\\
&\iff\beta\le\tfrac13\why{divide by $3>0$}
\end{align*}
With \(\beta\le1/3\), a liquidity abstention is at least as likely as a liquidity order in one direction.
Now repeat the argument for the nonzero orders with \(b_0\) in place of \(b\).
The zero order has \(b_0\le f_j(0)\le b_0+\lambda\), so its likelihood ratio lies between \(1/(1+\lambda/b_0)\) and \(1+\lambda/b_0\).
\begin{align*}
1+\frac{\lambda}{b_0}&\le1+\frac{\lambda}{b}\why{\(b_0\ge b>0\): divide \(\lambda\) by the larger number}\\
&=1+u\why{definition of \(u\)}
\end{align*}
So the same bound \(1+u\) holds for every order, and the lower bound \(1/(1+u)\) follows in the same way.

\begin{watchout}
Do not harmonize the domains of \(\beta\).
The slot is defined with \(0<\beta\le1/3\).
The posterior bound common to all orders, and the exclusion bound of Corollary~A.1, use \(\beta\le1/3\).
Proposition~2 states \(0<\beta<1/2\), and its proof needs bounds only at the nonzero orders (Appendix A.3).
Copy the domain of each result from its statement.
If \(\beta>1/3\), then \(b_0<b\).
A profile in which class \(H\) always abstains and class \(L\) never does then gives the zero order a likelihood ratio above \(1+u\).
\end{watchout}

\begin{checkyourself}
Show that \(b_0\ge b\) holds exactly when \(\beta\le1/3\).
Then take \(\beta>1/3\), the profile \(\sigma_H=\delta_0\) (class \(H\) always abstains), and \(\sigma_L(0)=0\).
Compute \(f_H(0)/f_L(0)\) and show that it exceeds \(1+u\).
\label{q:model-beta}
\end{checkyourself}

{\small\color{notegrey}Source: Section~2.3 of the paper, and Appendix A.1 and A.2 (Lemma A.1).}

%==================================================================
\section{Prices and the investor's profit}\label{sec:model-prices}

The competitive market makers see the order \(X\) and set the price
\[
P(x)=\E[V_T\mid X=x],
\]
where \(V_T\) is the terminal payoff of one target share.
Each share held at the sale receives \(V_T\), and market-maker inventory is no exception.

The rule has a reason.
Suppose the order is a purchase. The market maker sells one share at \(P(x)\) and gives up the payoff \(V_T\) of that share. Its profit is \(P(x)-V_T\).
Suppose the order is a sale. The market maker buys one share at \(P(x)\), and later receives \(V_T\). Its profit is \(V_T-P(x)\).

Competition drives the expected profit on each order to zero.
Given the order, \(P(x)\) is a fixed number, so the expected profit is \(\pm\bigl(P(x)-\E[V_T\mid X=x]\bigr)\).
This is zero exactly when \(P(x)\) equals the expectation of \(V_T\) given the order.
The paper uses the same rule after an abstention, and the challenger sees that price too.

\begin{watchout}
The terminal payoff \(V_T\) depends on whether the challenger prepares.
The challenger decides after it sees the price.
So the price depends on the payoff, and the payoff depends on the price.
The equilibrium conditions resolve this loop. Chapters~\ref{ch:beliefs} and~\ref{ch:trading} study it.
\end{watchout}

\subsection*{The profit of the investor}

Suppose the investor fills the slot with the order \(q\).
Its order is the whole order that the market makers see, so it trades at the price \(P(q)\).
Check the three cases.
\begin{itemize}
\item \(q=1\): the investor buys one unit at \(P(1)\) and receives \(V_T\). The profit is \(V_T-P(1)-k\).
\item \(q=-1\): the investor sells one unit at \(P(-1)\) and owes \(V_T\). The profit is \(P(-1)-V_T-k\).
\item \(q=0\): the investor does nothing. The profit is zero.
\end{itemize}
One formula covers the three cases:
\begin{equation}\label{eq:model-profit}
q\,\bigl[V_T-P(q)\bigr]-k\,|q|.
\end{equation}
For \(q=-1\) the formula gives \(-[V_T-P(-1)]-k=P(-1)-V_T-k\). For \(q=0\) it gives zero.

Three remarks on the symbols follow.
\begin{itemize}
\item The unit is a normalized small trading unit.
\item The cost \(k\) is a separate friction, such as an execution cost or a cost to carry the position. It is distinct from adverse-selection price impact, which the competitive price already produces.
\item The investor holds no net position, so a sale after bad news is a short sale. The cost \(k\) then includes the cost to borrow shares.
\end{itemize}

What this says: the investor gains only when its own value for the share differs from the price in the direction of its trade.
A purchase gains when the value is above the price, and a sale gains when the value is below it.
The price already contains the market makers' guess that the order came from the investor. The investor earns the residual advantage, and \(k\) takes some of it.
Chapter~\ref{ch:trading} asks when the residual covers \(k\).

{\small\color{notegrey}Source: Section~2.3 of the paper.}

%==================================================================
\section{The sequence of events}\label{sec:model-timing}

The sequence of events is fixed.
Figure~\ref{fig:model-timeline} shows the eight dates, and the list after it says what happens at each one.

\begin{figure}[tbp]
\centering
\begin{tikzpicture}[font=\small,>=Stealth]
\draw[->,thick] (-0.4,0) -- (13.6,0);
\foreach \x/\n/\lab/\pos in {
  0/1/{Seller announces the sale rule}/above,
  1.75/2/{Investor learns the class \(j\)}/below,
  3.5/3/{Slot draw selects the trader}/above,
  5.25/4/{Trader places the order \(X\)}/below,
  7.0/5/{Market makers set the price \(P\)}/above,
  8.75/6/{Challenger sees \(P\), learns \(C\), decides}/below,
  10.5/7/{Bidders learn values and bid}/above,
  12.25/8/{Ownership and payoffs are realized}/below}{
  \node[circle,draw,fill=white,inner sep=1.2pt,minimum size=14pt] (n\n) at (\x,0) {\(\n\)};
  \node[\pos=0.45cm of n\n,text width=2.9cm,align=center] {\lab};
}
\end{tikzpicture}
\caption{The timeline of the model. Time runs from left to right.}\label{fig:model-timeline}
\end{figure}

\begin{enumerate}
\item The seller announces the sale rule: a cash second-price auction with reserve \(p\).
\item The investor learns the class \(j\) of the challenger.
\item The slot draw selects the investor with probability \(\lambda\) and a liquidity trader otherwise.
\item The selected trader places its order.
\item The market makers see the order and set the price from it.
\item The challenger sees the price. After that, it learns its cost \(C\) in private and decides whether to prepare.
\item Prepared bidders learn their values and submit their bids.
\item Ownership and the financial payoffs are realized.
\end{enumerate}
Chapter~\ref{ch:research} adds a research stage before the investor learns the class.

The paper does not give a date at which nature draws the class, the values, and the cost.
The investor sees the class at date 2, so the class exists by then.
The challenger's value depends only on the class, and the cost and the incumbent's value are independent of everything else.
So the date of each draw does not matter. Only the date at which someone observes it matters.
This paragraph is the guide's own reading.

\subsection*{Who observes what}

Table~\ref{tab:model-info} lists what each party observes.
\begin{table}[tbp]
\centering\small
\caption{What each party observes. The columns are the class \(j\), the challenger's value \(\theta\), the incumbent's value \(R\), the cost \(C\), the order \(X\), and the price \(P\). A dash means that the paper does not say, and that no choice depends on it. The investor's entries ``if it trades'' are the guide's reading.}\label{tab:model-info}
\begin{tabular}{@{}lcccccc@{}}
\toprule
 & \(j\) & \(\theta\) & \(R\) & \(C\) & \(X\) & \(P\)\\
\midrule
Market makers & no & no & no & no & yes & they set it\\
Investor & yes & no & no & no & if it trades & if it trades\\
Challenger & no & after it pays \(C\) & no & yes & no & yes\\
Incumbent & -- & -- & yes & no & -- & --\\
\bottomrule
\end{tabular}
\end{table}

The paper states these entries as follows.
The market makers observe the order and nothing else.
The challenger observes the price and its own cost, but not the order, its class, its value, or the incumbent's value.
The investor observes the class but not the value.
The incumbent knows its own value.
The paper also says that the challenger learns its cost in private and that the incumbent's realized value stays private until it bids.
So the investor and the market makers observe neither \(C\) nor \(R\), and the incumbent does not observe \(C\).

Two entries are the guide's reading.
\begin{itemize}
\item \emph{The investor's order and price.} When the investor fills the slot, its own order is the order \(X\), and it trades at the price \(P(q)\), so it knows both. When a liquidity trader fills the slot, the investor makes no choice, so nothing depends on what it sees.
\item \emph{The incumbent.} Its bid equals its value whatever else it sees, so the paper needs to say nothing more about what it observes.
\end{itemize}

{\small\color{notegrey}Source: Sections~2.2 and 2.3 of the paper.}

%==================================================================
\section{The payoff objects}\label{sec:model-payoffs}

This section names the objects that the rest of the guide uses.
Chapter~\ref{ch:auction} derives their formulas. Here you learn what they mean.

\subsection*{Proceeds, gross profit, and the spread}

Suppose the challenger prepares and has value \(\theta\).
The auction then gives the target an expected amount \(t(\theta)\).
It gives the challenger an expected gross acquisition profit \(g(\theta)\).
Both expectations average over the incumbent's value \(R\), with truthful bids, and \(g\) is the profit before the cost \(C\).
For a value \(\theta\ge p\), the second-price rule splits the value between the target and the challenger, so \(t(\theta)+g(\theta)=\theta\) (Chapter~\ref{ch:auction}).

Average over the class.
Write \(\E_j\) for the expectation over \(\theta\sim G_j\), and define
\[
t_j=\E_j[t(\theta)],\qquad g_j=\E_j[g(\theta)],\qquad j\in\{H,L\}.
\]
The number \(t_j\) is the expected target proceeds after entry in class \(j\), and \(g_j\) is the expected gross acquisition profit in class \(j\).
The target-payoff spread is
\[
\DT=t_H-t_L.
\]
It tells how much more a class-\(H\) challenger adds to the target's proceeds than a class-\(L\) challenger.
A target share pays the proceeds of the sale, so the spread measures how much the value of a share depends on the class.

Bids are truthful, so these objects depend on the incumbent's distribution and the reserve, and on no belief.
When the incumbent's distribution matters, write \(g_j(F)\), \(\DT(F)\), and \(B_F\).
Along a path, write \(g_j(r)\), \(\DT(r)\), and \(B_r\) for the objects at \(F_r\).

\subsection*{Gross profit at a belief}

The challenger sees only the price \(\bar P\). It holds a belief \(\mu=\Pr(j=H\mid\bar P)\) at that price.
Its expected gross profit at the belief \(\mu\) is the function \(B(\mu)\):
\begin{align*}
B(\mu)&=\E\bigl[g(\theta)\mid\bar P\bigr]\why{definition}\\
&=\Pr(L\mid\bar P)\,\E\bigl[g(\theta)\mid L,\bar P\bigr]+\Pr(H\mid\bar P)\,\E\bigl[g(\theta)\mid H,\bar P\bigr]\why{total expectation}\\
&=(1-\mu)\,g_L+\mu\,g_H\why{given the class, \(\theta\) has law \(G_j\)}\\
&=g_L+\mu\,(g_H-g_L)\why{expand and collect}
\end{align*}
Here the third line uses the independence of the value and the order, given the class.
The price depends on the class and the slot draws only, so given the class it carries no news about \(\theta\).
Hence \(\E[g(\theta)\mid j,\bar P]=\E_j[g(\theta)]=g_j\).

What this says: the value depends on the price only through the class, so the profit at a belief is a weighted average of two numbers.
At \(\mu=0\) the profit is \(g_L\), and at \(\mu=1\) it is \(g_H\).
In between, it moves along a straight line. It rises with \(\mu\) exactly when \(g_H>g_L\).

\begin{checkyourself}
Name the independence assumption that makes \(B(\mu)\) linear in \(\mu\).
Suppose the price also carried news about \(\theta\) beyond the class.
Say which step of the derivation of \(B(\mu)\) fails, and what \(g_j\) would become.
\label{q:model-linear}
\end{checkyourself}

\subsection*{The gain from entry}

Define the gain in expected proceeds that entry brings in class \(j\), and its average at the belief \(\mu\) (Appendix A.2 of the paper):
\[
w_j=t_j-t_0,\qquad w(\mu)=(1-\mu)\,w_L+\mu\,w_H.
\]
Two identities connect these to the spread.
\begin{align*}
w_H-w_L&=(t_H-t_0)-(t_L-t_0)\why{definitions}\\
&=t_H-t_L=\DT\why{the terms $t_0$ cancel}\\[0.4em]
w(\mu)&=w_L+\mu\,(w_H-w_L)\why{expand $(1-\mu)w_L$}\\
&=t_L-t_0+\DT\,\mu\why{substitute $w_L$ and the first identity}
\end{align*}
Section~\ref{sec:model-equilibrium} uses \(w(\mu)\) in the price formula.

\subsection*{Entry}

Entry after the order \(x\), written \(e(x)\), is the probability that the challenger prepares after the price \(P(x)\).
Let \(\zeta(P(x))\) be the probability that a high-cost challenger prepares after that price.
Total probability over the cost gives:
\begin{align*}
e(x)&=\Pr(C=0)\cdot1+\Pr(C=c_H)\cdot\zeta(P(x))\why{total probability over $C$}\\
&=\rho+(1-\rho)\,\zeta(P(x))\why{$\Pr(C=0)=\rho$; the zero-cost challenger always prepares}
\end{align*}
The step uses the independence of the cost from the class and from the slot draws. So \(\Pr(C=0)=\rho\) given any class and any order.
Hence \(e(x)\) is also the probability of entry given the class and the order.
It follows that \(e(x)\ge\rho\) for every order, because \(\zeta\ge0\).
The floor \(\rho\) is part of entry. Entry is the probability that the challenger prepares, and a zero-cost challenger prepares.

Test the extremes.
If \(\zeta=0\), only the zero-cost challenger prepares, and \(e(x)=\rho\).
If \(\zeta=1\), every challenger prepares, and \(e(x)=1\).
If \(\rho\) is near one, entry is near one at every price.
If \(\rho\) is near zero, entry is close to \(\zeta\), and the high-cost challenger decides almost alone.

\begin{checkyourself}
Derive \(e(x)=\rho+(1-\rho)\zeta(P(x))\) from the cost distribution.
Show that \(e(x)\ge\rho\), and say what \(e(x)\) is when \(\zeta=0\) and when \(\zeta=1\).
Which independence assumption do you use?
\label{q:model-entry}
\end{checkyourself}

Entry in class \(j\), written \(e_j\), is the probability that the challenger prepares given the class \(j\).
Total entry \(\Ent\) is the probability that the challenger prepares, averaged over the class.
Both follow from total probability.
\begin{align*}
e_j&=\sum_x\Pr(X=x\mid j)\,\Pr(\text{prepares}\mid j,X=x)\why{total probability over the order}\\
&=\sum_xf_j(x)\,e(x)\why{definition of $f_j$ and $e$}\\[0.4em]
\Ent&=\pi\,e_H+(1-\pi)\,e_L\why{total probability over the class}
\end{align*}

\subsection*{Wins, ownership, and the takeover share}

A prepared challenger with value \(\theta\ge p\) wins when its value exceeds the incumbent's value.
Given \(\theta\), this has probability \(\Pr(R<\theta)=F(\theta)\), because \(F\) is continuous.
Average over the class:
\begin{equation}\label{eq:model-W}
W_j(F)=\E_j\bigl[F(\theta)\,\1\{\theta\ge p\}\bigr].
\end{equation}
Class \(H\) outbids the incumbent \(F\) when \(W_H(F)=1\). Then a prepared class-\(H\) challenger beats the incumbent for sure.
Under the class ranking the indicator is one, so \(W_j(F)=\E_j[F(\theta)]\).
The guide's own illustration: if \(G_H\) is a point mass at \(h\), then \(W_H(F)=F(h)\), and this is one exactly when the incumbent's value exceeds \(h\) with probability zero.

High-class ownership \(\Own_H\) is the probability that a challenger of class \(H\) acquires the target (Section~2.4 of the paper).
The proof of Proposition~2 makes the event precise: it is the event that the challenger is of class \(H\) and acquires the target (Appendix A.3, Step~4).
So \(\Own_H\) is a joint probability, and it includes the prior \(\pi\).
It is not the probability that the challenger acquires the target given class \(H\); that conditional probability is \(\Own_H/\pi\).

The challenger acquires the target exactly when it prepares and then wins the auction.
Say that the challenger would win when \(\theta\ge p\) and \(R<\theta\).
With truthful bids, a prepared challenger with such values acquires the target.
A tie \(R=\theta\) has probability zero, so it changes no probability below.

The next derivation gives a product formula.
\begin{derivation}[high-class ownership]
Fix class \(H\). Given the class, the event ``prepares'' depends on the cost and the slot draws only.
The event ``would win'' depends on \(\theta\) and \(R\) only.
These two groups are independent given the class (Section~\ref{sec:model-costs}).
\begin{align*}
\Own_H&=\Pr(j=H\text{ and prepares and would win})\why{definition, as in Appendix A.3}\\
&=\Pr(j=H)\,\Pr(\text{prepares and would win}\mid H)\why{multiplication rule}\\
&=\pi\,\Pr(\text{prepares}\mid H)\,\Pr(\text{would win}\mid H)\why{independence given the class}\\
&=\pi\,e_H\,\Pr(\theta\ge p,\ R<\theta\mid H)\why{definition of \(e_H\)}\\
&=\pi\,e_H\,\E_H\bigl[\1\{\theta\ge p\}\,\Pr(R<\theta\mid\theta)\bigr]\why{condition on \(\theta\), then average}\\
&=\pi\,e_H\,\E_H\bigl[\1\{\theta\ge p\}\,F(\theta)\bigr]\why{\(R\) independent of \(\theta\); \(F\) continuous}\\
&=\pi\,e_H\,W_H(F)\why{definition \eqref{eq:model-W}}
\end{align*}
\end{derivation}
In a market in which only the zero-cost challenger prepares, \(e_H=\rho\), and ownership is \(\pi\rho W_H(F)\).
Proposition~2 states ownership in this form against the weak and the very strong incumbent.

Two factors set ownership: entry in class \(H\), and the chance that class \(H\) wins.
A result about ownership therefore needs a statement about both factors.
A stronger incumbent can raise the first factor through the price and lower the second.
For this reason the clean rise of ownership in Proposition~2 asks that class \(H\) outbid the stronger of the incumbents that it compares, \(W_H(F_1)=1\) (Chapter~\ref{ch:ccc}).

The takeover share of class \(H\) is the probability that the challenger is of class \(H\), given that the challenger acquires the target.
Write \(\Own_L=(1-\pi)\,e_L\,W_L(F)\) for the same product in class \(L\), by the same steps.
The challenger is of class \(H\) or of class \(L\), and these cases are disjoint, so \(\Pr(\text{acquires})=\Own_H+\Own_L\).
Then
\[
\Pr(j=H\mid\text{acquires})=\frac{\Pr(j=H,\text{acquires})}{\Pr(\text{acquires})}=\frac{\Own_H}{\Own_H+\Own_L}\whyinline{definition of conditional probability}
\]

{\small\color{notegrey}Source: Section~2.4 of the paper, and Appendix A.2 for \(w_j\). Appendix A.3 (Step~4 of the proof of Proposition~2) gives the independence argument and the product formula for ownership.}

%==================================================================
\section{The equilibrium}\label{sec:model-equilibrium}

An equilibrium lists what each agent does and believes.
Nobody gains by a change of its own action, the market makers' prices match the expected payoff, and beliefs follow Bayes' rule.
This section states the definition, reads it condition by condition, and derives the investor's objective.

\subsection*{The definition}

\begin{paperresult}{Definition A.1 (Equilibrium)}
Write \(\eta(-1)=\eta(1)=\beta\) and \(\eta(0)=1-2\beta\).
An equilibrium is a tuple \(\mathcal{E}=(\sigma_H,\sigma_L,\mu,P,\bar\mu,\zeta)\), together with the bids of the prepared bidders.
Here \(\sigma_j\) is a probability distribution on \(\{-1,0,1\}\), the order of the informed investor in class \(j\), and mixed and wrong-signed orders are allowed;
\(\mu(x)=\Pr(j=H\mid X=x)\) is the market makers' posterior;
\(P(x)\) is the price after each order;
\(\bar\mu(\bar P)=\Pr(j=H\mid P(X)=\bar P)\) is the challenger's belief at each price \(\bar P\) in the image of \(P\);
and \(\zeta(\bar P)\in[0,1]\) is the probability that a high-cost challenger prepares after the price \(\bar P\).
A zero-cost challenger prepares after every price.
For each order \(x\), \(f_j(x)=(1-\lambda)\eta(x)+\lambda\sigma_j(x)\) and \(e(x)=\rho+(1-\rho)\zeta(P(x))\).
The tuple satisfies six conditions.
\begin{enumerate}[label=(E\arabic*),leftmargin=3em]
\item Bayes consistency. \(\mu(x)=\pi f_H(x)/[\pi f_H(x)+(1-\pi)f_L(x)]\) for every order \(x\).
\item Challenger beliefs. For every price \(\bar P\) in the image, \(\bar\mu(\bar P)\) is the posterior of class \(H\) given the event \(\{x:P(x)=\bar P\}\), computed from \(f_H\) and \(f_L\).
\item Entry optimality. For every price \(\bar P\) in the image, \(\zeta(\bar P)=1\) if \(B(\bar\mu(\bar P))>c_H\) and \(\zeta(\bar P)=0\) if \(B(\bar\mu(\bar P))<c_H\). At equality, the challenger's tie rule applies.
\item Competitive prices. \(P(x)=\E[V_T\mid X=x]\) for every order \(x\).
\item Investor optimality. For each class \(j\), \(\sigma_j(q)>0\) only if \(q\) maximizes \(U_j(q)=q\,[t_0+e(q)(t_j-t_0)-P(q)]-k|q|\) over \(q\in\{-1,0,1\}\), with \(P(\cdot)\) and \(e(\cdot)\) held at their values in \(\mathcal{E}\).
\item Truthful bids. Each prepared bidder whose value meets the reserve bids its value.
\end{enumerate}
\end{paperresult}

\subsection*{The six conditions in words}

\begin{enumerate}[label=(E\arabic*),leftmargin=3em]
\item \emph{Bayes consistency.} The market makers apply Bayes' rule after each order. The posterior is the share of class \(H\) in the probability of that order, with each class weighted by its prior.
\item \emph{Challenger beliefs.} The challenger sees the price, not the order. Two orders can share a price. The challenger then pools them, and its belief is the posterior given the set of orders with that price.
\item \emph{Entry optimality.} A high-cost challenger compares its expected gross profit \(B\) at its belief with the cost \(c_H\). It prepares when \(B>c_H\), stays out when \(B<c_H\), and follows a tie rule when \(B=c_H\). When \(g_H>g_L\), the function \(B\) rises with the belief, and the rule is a threshold rule. The belief at which \(B(\mu)=c_H\) is the preparation threshold \(\tau\) (Chapter~\ref{ch:beliefs}).
\item \emph{Competitive prices.} The price after each order is the expected terminal payoff of a share, given the order.
\item \emph{Investor optimality.} In each class, the investor places an order with positive probability only if that order gives the largest payoff \(U_j\). It takes the prices and the entry rule as given.
\item \emph{Truthful bids.} Every prepared bidder whose value meets the reserve bids its value. Chapter~\ref{ch:auction} shows that this is weakly dominant, so (E6) holds in every tuple.
\end{enumerate}

\subsection*{Why every order has positive probability in both classes}

By \eqref{eq:model-orderprob}, the probability of an order obeys
\[
f_j(x)=(1-\lambda)\eta(x)+\lambda\sigma_j(x)\ \ge\ (1-\lambda)\eta(x)>0.
\]
The first inequality holds because \(\lambda\sigma_j(x)\ge0\).
The second holds because \(\lambda<1\) and \(\eta(x)>0\) for \(0<\beta<1/2\).
This holds under every investor strategy.
So every order has positive probability in both classes, Bayes' rule defines \(\mu(x)\) after every order, and \(0<\mu(x)<1\) because \(0<\pi<1\).
The challenger sees only prices that some order produces, so its beliefs at those prices also come from Bayes' rule.
No off-path belief appears anywhere in the definition.

The unit menu \(\{-1,0,1\}\) keeps every order on path.
A menu with continuous order sizes would need a convention for the sizes that no class uses (Edmans, Goldstein, and Jiang, 2015).

\subsection*{The investor's objective}

The next derivation finds the expected terminal payoff of a share given the class \(j\) and the order \(q\).

\begin{derivation}[the expected terminal payoff given the class and the order]
Let \(A\) be the event that the challenger prepares.
\begin{enumerate}[leftmargin=2em]
\item The challenger's choice depends on the price \(P(x)\) and the cost \(C\) only. The cost is independent of the class and the slot draws, so \(\Pr(A\mid j,X=x)=\rho+(1-\rho)\zeta(P(x))=e(x)\) (Section~\ref{sec:model-payoffs}).
\item Given the class, the pair \((\theta,R)\) is independent of the cost and of the slot draws. The order depends on the class and the slot draws only. So, given \(j\) and \(X=x\), the event \(A\) is independent of \((\theta,R)\).
\item On \(A\), bids are truthful by (E6). By step 2, \(\theta\) still has law \(G_j\), so the expected proceeds are \(\E_j[t(\theta)]=t_j\).
\item On the complement of \(A\), the proceeds are \(p\,\1\{R\ge p\}\). By step 2, their expectation is \(t_0\).
\end{enumerate}
Combine the two cases:
\begin{align*}
\E[V_T\mid j,X=x]&=e(x)\,t_j+\bigl[1-e(x)\bigr]\,t_0\why{total expectation over $A$ and its complement}\\
&=t_0+e(x)\,(t_j-t_0)\why{collect the terms in \(t_0\)}
\end{align*}
The same steps hold if you also condition on the event that the investor fills the slot. That event is a slot draw, so step 2 still applies.
\end{derivation}

Next, take the expected value of the profit \eqref{eq:model-profit} given the class \(j\) and the order \(q\).
The numbers \(q\), \(P(q)\), and \(k|q|\) are fixed once the order is fixed.
\begin{align*}
U_j(q)&=\E\Bigl[q\,\bigl(V_T-P(q)\bigr)-k|q|\ \Big|\ j,X=q\Bigr]\why{expected profit of the order $q$}\\
&=q\,\bigl(\E[V_T\mid j,X=q]-P(q)\bigr)-k|q|\why{linearity of expectation; $q$ and $P(q)$ are fixed}\\
&=q\,\bigl[t_0+e(q)\,(t_j-t_0)-P(q)\bigr]-k|q|\why{substitute the derivation above}
\end{align*}
This is the objective in (E5).

What this says: the investor values a share at \(t_0+e(q)(t_j-t_0)\), because it knows the class \(j\).
The market makers value the share at \(P(q)\), because they know only the order.
The investor trades on the difference and pays \(k\).

The same independence gives the price formula.
Average the expected payoff over the class, with the weights \(\mu(x)\) and \(1-\mu(x)\) from (E1), and use \(w(\mu)\) from Section~\ref{sec:model-payoffs}:
\begin{align*}
P(x)&=(1-\mu(x))\,\E[V_T\mid L,X=x]+\mu(x)\,\E[V_T\mid H,X=x]\why{(E4), total expectation}\\
&=(1-\mu(x))\,\bigl[t_0+e(x)\,w_L\bigr]+\mu(x)\,\bigl[t_0+e(x)\,w_H\bigr]\why{terminal payoff; \(w_j=t_j-t_0\)}\\
&=t_0+e(x)\,\bigl[(1-\mu(x))\,w_L+\mu(x)\,w_H\bigr]\why{the weights add to one}\\
&=t_0+e(x)\,w(\mu(x))\why{definition of $w(\mu)$}\\
&=t_0+e(x)\,\bigl[t_L-t_0+\DT\,\mu(x)\bigr]\why{identity for $w(\mu)$}
\end{align*}
In the third line, \((1-\mu(x))\,t_0+\mu(x)\,t_0=t_0\), and \(e(x)\) multiplies both of the other terms, so it comes out of the bracket.
Display the result:
\begin{equation}\label{eq:model-price}
P(x)=t_0+e(x)\,\bigl[t_L-t_0+\DT\,\mu(x)\bigr].
\end{equation}

\begin{watchout}
The price is affine in the posterior \(\mu(x)\) at a fixed level of entry \(e(x)\).
Entry is not fixed in equilibrium: by (E3), it moves with the belief.
With a zero spread \(\DT\), the posterior drops out of \eqref{eq:model-price} at a given level of entry, so the price does not move with the posterior.
The price reveals the market makers' posterior only when the spread is positive and the entry floor \(\rho\) is positive.
Then entry is at least \(\rho>0\) after every order, and a higher posterior gives a strictly higher price (Section~3.2 of the paper).
Lemma~A.2 states the full conditions, and Chapter~\ref{ch:beliefs} proves it.
\end{watchout}

\subsection*{Why the factor \texorpdfstring{\(\lambda\)}{lambda} does not change the best order}

Before the slot draw, the investor of class \(j\) earns \(U_j(q)\) with probability \(\lambda\sigma_j(q)\) for each order \(q\).
It earns zero if a liquidity trader fills the slot.
Its ex ante payoff is
\[
\lambda\sum_q\sigma_j(q)\,U_j(q)+(1-\lambda)\cdot0.
\]
The sum is a weighted average of the three numbers \(U_j(q)\), with the weights \(\sigma_j(q)\).
A weighted average is at most the largest of the numbers.
It equals that number exactly when the weight sits on orders that attain it.
The positive factor \(\lambda\) multiplies the whole expression.
So the investor maximizes its payoff when it puts weight only on orders that maximize \(U_j\), and this is condition (E5).

{\small\color{notegrey}Source: Section~2.4 of the paper, Definition A.1, and Lemma A.2 (Step 0) in Appendix A.2.}

%==================================================================
\section{Ties}\label{sec:model-ties}

An indifference can arise in two places: the investor can have two best orders, and a high-cost challenger can have \(B=c_H\) exactly.
The paper calls two rules the default tie rules:
\begin{itemize}
\item an indifferent investor may mix over its best orders;
\item an indifferent high-cost challenger prepares.
\end{itemize}
The paper also considers two other rules:
\begin{itemize}
\item an indifferent investor abstains;
\item an indifferent high-cost challenger may prepare with any probability.
\end{itemize}
Each combination of a rule for the investor and a rule for the challenger is a tie rule of the paper, so there are four tie rules.
In Definition A.1, the default rules apply at the two indifferences, (E3) for the challenger and (E5) for the investor.

A mix over two orders needs indifference.
Condition (E5) puts positive weight only on orders that maximize \(U_j\).
The investor therefore mixes only when two or more orders give the same payoff.
The tie rules decide what happens at these points.

Where another tie rule changes a conclusion, the appendix states the change beside the result.
A table in Appendix~A.4 collects the changes for coexistence in the example at the prior \(1/2\).
The paper states that Proposition~2 holds under all four tie rules.
Propositions~3 and~4 take the default rules.
Each later chapter names the tie rule of each result it presents.

\subsection*{A unilateral deviation and a comparison across economies}

These two ideas answer different questions.

\begin{itemize}
\item A unilateral deviation by the investor changes only its own order. It faces the candidate equilibrium's price function and entry rule. The deviation payoff is therefore \(U_j(q)\) read at the candidate's \(P(q)\) and \(e(q)\).
\item A comparison across economies, such as across incumbent strengths, solves for the equilibrium again in each economy. The prices, beliefs, entry, and orders can all change from one economy to the next.
\end{itemize}

The first idea checks whether a candidate is an equilibrium.
The second idea asks how the equilibrium outcome moves when a parameter moves.

\begin{watchout}
A deviation faces fixed prices and entry.
A comparison across incumbent strengths solves the equilibrium again in each economy.
The paper's results on the effect of a stronger incumbent are comparisons of this second kind.
\end{watchout}

{\small\color{notegrey}Source: Section~2.4 of the paper, and the conventions of Appendix A.2.}

%==================================================================
\section{The money normalization}\label{sec:model-money}

The paper makes two statements about the unit of money.

The first concerns the target's known standalone value, which the paper normalizes to zero.
A common standalone component added to every ownership outcome shifts prices and payoffs by a constant, and it leaves incentives unchanged.
In the investor's objective the reason is short: the shift raises \(V_T\) and \(P(q)\) by the same constant, so the difference \(V_T-P(q)\) does not change.

The second is the rescale.
A common positive rescale of every value, the reserve, and every preparation, trade and research cost rescales prices and payoffs by the same factor.
Orders, beliefs, entry, and ownership do not change.
No normalization of money is needed for the general statements.

Let \(s>0\) be the common factor.
The values \(R\) and \(\theta\), the reserve \(p\), and the costs \(c_H\) and \(k\) are all multiplied by \(s\).
The probabilities \(\pi\), \(\rho\), \(\lambda\), and \(\beta\) do not change.
Take an equilibrium of the original economy, and build the rescaled tuple from it.
The rescaled tuple keeps the orders \(\sigma_j\) and the posteriors \(\mu\), and it multiplies every price by \(s\).
At the rescaled price \(s\bar P\), it gives the challenger the old belief \(\bar\mu(\bar P)\) and the old preparation probability \(\zeta(\bar P)\).
The list below checks that this tuple is an equilibrium of the rescaled economy.

\begin{itemize}
\item \emph{Orders and beliefs.} The order probabilities \(f_j(x)\) use only \(\lambda\), \(\beta\), and \(\sigma_j\). They contain no money. Bayes' rule uses only \(\pi\) and \(f_j\). So with the same orders, the beliefs are the same, and two orders share a rescaled price exactly when they shared the old price.
\item \emph{Proceeds.} The incumbent's rescaled value \(sR\) has distribution \(F_s(y)=F(y/s)\). The rescaled reserve is \(sp\). Then \(F_s(sp)=F(p)\), so the rescaled no-entry proceeds are \(sp\,[1-F(p)]=s\,t_0\). Every other payment in the auction is the reserve, a bid, or a value. So \(t(\theta)\), \(g(\theta)\), \(t_j\), \(g_j\), and \(\DT\) are all multiplied by \(s\), and so are \(V_T\) and each price \(\E[V_T\mid X=x]\).
\item \emph{Entry.} The expected gross profit \(B(\mu)\) is multiplied by \(s\), and the cost \(c_H\) is multiplied by \(s\). The comparison \(B(\mu)\gtrless c_H\) has the same answer. The entry rule is the same.
\item \emph{Ownership.} Who wins depends on comparisons among \(\theta\), \(R\), and \(p\). When you multiply all three by \(s>0\), every comparison keeps its answer. So the chance \(W_j\) that a prepared class-\(j\) challenger wins, and ownership, do not change.
\item \emph{The investor.} The order \(q\) is in \(\{-1,0,1\}\), so \(s\) does not rescale it. Each of \(t_0\), \(t_j\), \(P(q)\), and \(k\) is multiplied by \(s\), so \(U_j(q)\) is multiplied by \(s\). The ranking of the orders does not change, and the best order does not change.
\end{itemize}

These facts show that each condition (E1) to (E6) holds in the rescaled economy for the rescaled tuple.
Conditions (E1), (E2), (E3), and (E6) hold for the same reasons as before. In conditions (E4) and (E5), each side is multiplied by \(s\).

\begin{checkyourself}
Double all money, so \(s=2\).
Show that the no-entry proceeds \(t_0\) double.
Show that \(U_j(q)\) doubles for each order \(q\), and say why the best order stays the same.
\label{q:model-scale}
\end{checkyourself}

{\small\color{notegrey}Source: Section~2.2 of the paper.}

%==================================================================
\begin{chaptersummary}
\begin{itemize}
\item The model has six kinds of agent: the target with its seller, the incumbent, the challenger, the informed investor, the liquidity trader, and the market makers.
\item The class ranking says that class \(H\) dominates class \(L\) and that every value meets the reserve.
\item The slot gives each order the probability \(f_j(x)=(1-\lambda)\eta(x)+\lambda\sigma_j(x)\), so one order carries a likelihood ratio of at most \(1+u\); for the zero order this uses \(\beta\le1/3\).
\item The auction without entry pays the target \(t_0=p[1-F(p)]\), and entry \(e(x)\) never falls below \(\rho\).
\item An equilibrium has six conditions, every order has positive probability in both classes, and the investor's objective is \(U_j(q)=q[t_0+e(q)(t_j-t_0)-P(q)]-k|q|\).
\item A deviation faces fixed prices and entry, and a comparison across economies solves the equilibrium again.
\item A common rescale of money changes no order, belief, entry, or ownership.
\end{itemize}
\end{chaptersummary}

%==================================================================
\answersection

\answer{\ref{q:model-t0}}{%
Part one.
If \(p\ge\bar r\), then \(F(p)=1\), because \(F\) is extended by unity above \(\bar r\) and \(F(\bar r)=1\).
Then \(t_0=p[1-F(p)]=p\cdot0=0\).
In words: the incumbent's value is below the reserve for sure, so no sale occurs and the target receives nothing.

Part two.
Since \(F(p)\) is a probability, \(0\le F(p)\le1\), so \(0\le1-F(p)\le1\).
Multiply by \(p>0\) and the inequality direction stays the same: \(0\le p[1-F(p)]\le p\).

Part three.
In the example, \(F(p)=p/r\) for \(0<p<r\).
Then \(t_0=p(1-p/r)=p(r-p)/r\), a product of \(p\) and \(1-p/r\).
The first factor is zero as \(p\) tends to zero, and the second factor is zero at \(p=r\).
So \(t_0\) is zero at both ends and positive between them.}

\answer{\ref{q:model-beta}}{%
Compare \(b_0\) and \(b\):
\begin{align*}
b_0\ge b&\iff(1-\lambda)(1-2\beta)\ge(1-\lambda)\beta\why{definitions}\\
&\iff1-2\beta\ge\beta\why{divide by $1-\lambda>0$}\\
&\iff1\ge3\beta\why{add $2\beta$}\\
&\iff\beta\le\tfrac13\why{divide by $3>0$}
\end{align*}
Let \(\beta>1/3\), so \(b_0<b\). Also \(\beta<1/2\), so \(b_0>0\).
With \(\sigma_H=\delta_0\), \(f_H(0)=b_0+\lambda\cdot1\).
With \(\sigma_L(0)=0\), \(f_L(0)=b_0+\lambda\cdot0=b_0\).
The ratio is
\[
\frac{f_H(0)}{f_L(0)}=\frac{b_0+\lambda}{b_0}=1+\frac{\lambda}{b_0}.
\]
Divide \(\lambda\) by the smaller positive number \(b_0<b\): then \(\lambda/b_0>\lambda/b=u\).
So the ratio exceeds \(1+u\), and the bound \(1+u\) fails for the zero order.
}

\answer{\ref{q:model-linear}}{%
The assumption is that, given the class, the challenger's value \(\theta\) is independent of the order, and so of the price.
It lets you replace \(\E[g(\theta)\mid j,\bar P]\) by \(\E_j[g(\theta)]=g_j\) in the third line of the derivation.
Suppose the price also carried news about \(\theta\) beyond the class.
The third line would then contain \(g_j(\bar P)=\E[g(\theta)\mid j,\bar P]\), a number that changes with the price.
The expected profit would be \((1-\mu)g_L(\bar P)+\mu g_H(\bar P)\).
It would not be a straight line in \(\mu\) with the fixed ends \(g_L\) and \(g_H\).
}

\answer{\ref{q:model-entry}}{%
The challenger has the cost \(0\) with probability \(\rho\) and the cost \(c_H\) with probability \(1-\rho\).
A zero-cost challenger prepares for sure. A high-cost challenger prepares with probability \(\zeta(P(x))\).
Total probability over the cost gives
\[
e(x)=\rho\cdot1+(1-\rho)\,\zeta(P(x)).
\]
The step uses the independence of the cost from the class and the slot draws: the probability \(\Pr(C=0)=\rho\) holds after every order and in every class.

Since \(\zeta\ge0\) and \(1-\rho>0\), the second term is nonnegative, so \(e(x)\ge\rho\).
If \(\zeta=0\), then \(e(x)=\rho\): only the zero-cost challenger prepares.
If \(\zeta=1\), then \(e(x)=\rho+(1-\rho)=1\): every challenger prepares.
}

\answer{\ref{q:model-scale}}{%
Let \(s=2\).
The incumbent's value becomes \(2R\), with distribution \(F_2(y)=\Pr(2R\le y)=\Pr(R\le y/2)=F(y/2)\).
The reserve becomes \(2p\), and \(F_2(2p)=F(p)\).
The no-entry proceeds of the doubled economy are \(2p[1-F_2(2p)]=2p[1-F(p)]=2t_0\).

The proceeds \(t_j\), the price \(P(q)\), and the cost \(k\) double in the same way, and the order \(q\in\{-1,0,1\}\) does not change.
The entry \(e(q)\) does not change, because it depends on the cost comparison \(B(\mu)\gtrless c_H\), and both sides double.
Write \(\tilde U_j\) for the objective of the doubled economy. Then
\[
\tilde U_j(q)=q\bigl[2t_0+e(q)(2t_j-2t_0)-2P(q)\bigr]-2k|q|=2\,U_j(q).
\]
Multiplication by \(2\) keeps the ranking of the three numbers \(U_j(-1)\), \(U_j(0)\), \(U_j(1)\): \(2U_j(q)\ge2U_j(q')\) exactly when \(U_j(q)\ge U_j(q')\).
So the best order, and every tie, stays the same.
}
